\paragraph{Why eventization helps beyond compression.}
Post-hoc compression reduces what is \emph{stored}; acquisition-side
eventization reduces what is \emph{sensed and transmitted}, which is
the binding constraint for battery-powered and wireless SHM nodes
\cite{henkmann2025tinyml}. Our Pareto comparison (E2) quantifies
whether the information lost by event sampling exceeds that lost by
transform coding at the same rate.

\paragraph{Event structure as a physical feature.}
The coherence--burstiness separation rests on physics: temperature is a
global, slow boundary condition, whereas damage is a local, persistent
scatterer. This makes the discriminator interpretable and, unlike
learned compensators, free of a training distribution over temperature.

\paragraph{Limitations.}
(i) The level-A residual quality still depends on compensation; level B
provides an independent channel but at coarser time resolution.
(ii) The OGW damage cases are reversible surface-mounted discs, not
grown defects; conclusions about small real cracks require caution.
(iii) The CFRP plate is quasi-isotropic, so directional (anisotropy)
claims are stated conservatively. (iv) Eight of thirteen damage stages
of the long-term dataset are confounded with sensor-drift events; we
report those separately rather than merging them into headline metrics.

\paragraph{Threats to validity.}
Synthetic-to-real gap: pipeline sanity checks used a simplified
dispersion/temperature model; all headline numbers are from the two
public real datasets.
